%% routeid --- workshop/short-paper draft.
%% Compile on Overleaf (pdfLaTeX) with acmart.cls and refs.bib.
\documentclass[sigconf,nonacm]{acmart}
\settopmatter{printacmref=false}
\usepackage{booktabs}
\usepackage{graphicx}
\usepackage{url}
\setlength{\emergencystretch}{3em}

\begin{document}

\title{RouteID: Memory-Mapped, Locality-Ordered Routing for Bounded-Radius Routes on Indonesian OpenStreetMap Data}

\author{Anonymous}
\affiliation{\institution{Preprint}\country{}}
\renewcommand{\shortauthors}{Anonymous}

\begin{abstract}
Routing engines for OpenStreetMap (OSM) are engineered for planet-scale,
long-distance queries and trade memory for hierarchy. Many production workloads
are the opposite: bounded-radius queries, typically below 100\,km, issued from a
mobile app that needs a polyline and a distance estimate, and billed per request
when a hosted API is used. We ask what such a workload actually requires and
present \emph{RouteID}, a Go engine (6.2\,kLOC for the routing core and HTTP
server) that compiles OSM extracts into one memory-mapped binary graph and
answers queries with edge-based, heuristic-inflated A*. On the Java extract
(22.71\,M nodes, 46.68\,M edges) the evaluation graph is 693\,MB and is built in
28\,s; the process starts at 17.3\,MB resident and answers 332 filter-matched
pairs at p50 38.7\,ms and p95 674\,ms, versus OSRM MLD at 1.5\,ms and 3.7\,ms.
RouteID keeps a 7.3$\times$ smaller artefact and a 140$\times$ smaller idle
resident set, but it is slower; we do not claim otherwise. Two inexpensive
techniques are evaluated in isolation: Hilbert-curve node renumbering lowers p50
by 7--20\,\%, and a 1.5 heuristic inflation is 28\,\% faster while reproducing the
optimal route length on 526 of 526 queries. A single national graph (54.89\,M
nodes, 112.34\,M edges, 1.621\,GB) is built in 11\,m40\,s including download.
On the distance comparison, the median relative difference is 0.23\,\% but p95 is
14.49\,\% and the maximum is 70.44\,\%. We did not collect an independent
reference, so no accuracy claim is made. At a 100\,km cap, 481 of 600 pairs
succeed and p95 wall-clock latency is 0.94--1.32\,s; the tail, and 8 budget
failures, are the price of omitting a hierarchy.
\end{abstract}

\ccsdesc[500]{Information systems~Geographic information systems}
\ccsdesc[500]{Theory of computation~Shortest paths}
\keywords{OpenStreetMap, route planning, memory-mapped graphs, spatial locality,
weighted A*, Indonesia}

\maketitle

\section{Introduction}
Indonesia is an archipelagic country whose road data in OpenStreetMap is uneven:
dense and well mapped in Java and Bali, sparse in Papua and Maluku. A large class
of consumer applications---ride hailing, delivery, field service---does not need
continental routing. It needs, for two coordinates supplied by the user, a road
polyline within a bounded radius that can be drawn on a basemap, plus a distance
estimate and a coarse ETA.

Hosted directions APIs are the default for that workload, but they are priced per
request. Self-hosted engines such as OSRM, GraphHopper, and Valhalla are free per
request, but are designed and tuned for long queries; their memory footprint
(roughly a gigabyte or more for a national extract) dominates the cost of a cheap
virtual machine. Our question is narrow: \emph{for bounded-radius workloads, what
does a purpose-built engine have to do, and where does it lose to a
general-purpose one?} We report an engineering case study, not a new shortest-path
algorithm and not an accuracy result.

Our contributions are:
\begin{itemize}
  \item an extraction and serving pipeline whose graph is a single memory-mapped
        file, with one artefact carrying car and motorcycle profiles, and whose
        resident set scales with touched pages rather than with the data set;
  \item an isolated evaluation of two inexpensive choices: Hilbert-curve node
        renumbering and A* heuristic inflation;
  \item an edge-based search that enforces OSM turn restrictions with a sparse
        transition table, plus segment snapping that removes most attachment
        error;
  \item a controlled comparison against OSRM MLD on the same machine and data,
        including the negative result at the distance tail;
  \item a national build and a 100\,km-cap capacity measurement, with explicit
        limits: no independent ground truth, no traffic, and no full-national
        latency benchmark.
\end{itemize}

\section{Workload and requirements}
\begin{table}[t]
\small
\begin{tabular}{@{}ll@{}}
\toprule
R1 & routes $\le 100$\,km; most queries $\le 50$\,km; car and motorcycle \\
R2 & output: encoded polyline, distance, coarse ETA \\
R3 & no per-request fee; commodity VPS with 4\,GB RAM \\
R4 & memory-mapped graph; restart without a warm-up cost \\
R5 & Indonesian data: ferries, one-way systems, dense turn restrictions \\
R6 & offline-capable build; monthly OSM refresh \\
\bottomrule
\end{tabular}
\caption{Requirements that shaped the design. They motivate a memory-first
engine; they do not prove that a hierarchy is the wrong choice.}
\label{tab:req}
\end{table}

Table~\ref{tab:req} is deliberately modest. We do not require traffic, worldwide
coverage, matrix queries, map matching, or turn-by-turn instructions. The output
is a line to draw and a scalar distance; ETA is static-speed and explicitly not
authoritative. The 100\,km cap is a product decision: beyond it the absence of a
hierarchy dominates, as the tail measurements show.

\section{Design}
\subsection{Extraction and profiles}
The builder reads an \texttt{.osm.pbf} twice (ways, then nodes) and keeps
car-routable and motorcycle-routable ways. It applies per-class default speeds
tuned for Indonesian conditions (residential 25\,km/h, service 15\,km/h, primary
60\,km/h), \texttt{maxspeed} overrides, \texttt{access} and
\texttt{motor\_vehicle} filtering, \texttt{oneway}, and
\texttt{junction=roundabout}. Every OSM way node becomes a graph node, so stored
geometry is the OSM geometry and no separate shape-point table is needed.
\texttt{route=ferry} ways are kept at 20\,km/h. Format v5 stores both profiles in
one file: car classes and motorcycle classes share the topology, while
motorcycle-specific access is a separate class byte.

The fallback speed table was aligned with OSRM's \texttt{car.lua} because an
earlier mismatch was the largest single source of route-distance deviation. The
alignment does not make the engines equivalent: OSRM minimises its own
\texttt{weight}, while RouteID minimises drive time plus a turn penalty.

\subsection{Binary format v5}
The graph is one file: a 320-byte header followed by 14 aligned sections. The
header carries magic, version, counts, grid parameters, the bounding box, section
offsets, profile metadata, PBF SHA-256, and build time. Nodes are 12\,bytes:
fixed-point latitude and longitude at $10^{-7}$\,deg (8\,bytes) and the first
outgoing-edge index (4\,bytes). Edges are 8\,bytes: destination node (4\,bytes),
length in units of 4\,m stored as \texttt{uint16} (2\,bytes), car class
(1\,byte), and motorcycle class (1\,byte). A class byte is speed in km/h
multiplied by two; zero means no access. Travel time is \emph{computed}, not
stored: one length unit at class $c$ costs $288/c$ deciseconds. There is no
interior-geometry section and no time array; the CSR adjacency is
compressed-sparse-row, and a uniform grid over nodes supports snapping.

Two sparse tables in the header tail hold forbidden and permitted transition
pairs, separately for car and motorcycle. A workspace key uses 28 bits for an
edge identifier, so a graph with more than $2^{28}$ edges is rejected at build
time; the national graph has 112.34\,M and is well below that limit.

\subsection{Node ordering}
Node order is the layout knob that matters. We sort nodes along a Hilbert curve
over a $2^{16}\times2^{16}$ grid of the graph's bounding box. Because the A*
frontier explores a spatially coherent region, a Hilbert numbering turns that
region into a contiguous byte range, which the page cache can serve without
thrashing. We keep OSM-id order behind a flag (\texttt{-order osm}) for A/B
measurement. The ordering does not change the search or its expansion count; it
changes which pages are touched.

\subsection{Search}
The search state is the \emph{edge} just traversed, not the node. This costs
nothing in the data layout---successors of state $e$ are the CSR range of its head
node---and it is what makes turn restrictions expressible: an OSM restriction
constrains the pair (arrival edge, departure edge). Forbidden pairs are a sorted
array of packed keys; \texttt{only\_*} rules are a sorted array of permitted
pairs. A lookup is one binary search over a few thousand entries, and a route
through unrestricted territory pays a single key comparison per transition.

We run A* with an admissible geometric bound
\[
  h(e) = \varepsilon \cdot \frac{\|\mathrm{head}(e),t\|}{v_{\max}},
\]
with a 1\,\% lower-bound margin, where $v_{\max}=130$\,km/h and
$\varepsilon\ge1$. Costs are travel times plus a turn cost, so the returned route
is the fastest path under static speeds. This matches
how hosted APIs choose routes; minimising distance instead produces visibly
different polylines. The turn cost is read from the change of heading between
arrival and departure edge---two node coordinates per expansion, no new graph
field---bucketed at 22.5$^\circ$/67.5$^\circ$/135$^\circ$ into 0/2/7.5/20\,s, and
shipped at a quarter of that weight. That setting leaves the E1 distance spread
unchanged but removes the worst detours (\texttt{paper/RESULTS.md}\,\S15).

The open set is a binary heap. Distances and parents live in a generation-stamped
open-addressing table that grows at 70\,\% load, so no per-query clearing is needed
and a query allocates nothing. Routers hand out pre-allocated workspaces from a
shared pool, which bounds peak memory at the cost of reuse. A query is bounded at
$4\cdot10^6$ expansions and by a budget (60\,\% of the HTTP timeout, 6\,s by
default). If the corridor exceeds the bound, the search retries with a greedier
heuristic before reporting \texttt{budget}; both bounds protect tail latency at
the cost of some long routes.

\subsection{Snapping and output}
Queries are snapped to the nearest \emph{segment}: the grid (cell
$0.02^{\circ}$, about 2.2\,km) is expanded ring by ring, and for every node found
the outgoing segments are projected onto the local tangent plane
(equirectangular, sub-metre at city scale). The projection yields an anchor node
(the closer end of the segment) and an along-distance. After routing between
anchors, geometry is corrected without mutating the graph: if the path leaves the
anchor \emph{along} the snapped segment, the covered part is subtracted from
distance and time and the anchor is dropped from the geometry (\emph{trim});
otherwise the offset is added as a spur. Ring expansion stops once no segment can
beat the best projection, with 500\,m of slack for segments longer than their
node-to-node distance. The server rejects an endpoint with no road within 3\,km
with reason \texttt{snap}.

Output is a precision-5 encoded polyline (about 1.1\,m), optionally
Douglas--Peucker simplified. The projected points are returned as
\texttt{snap\_from} and \texttt{snap\_to}; \texttt{pin=1} restores the older
behaviour that appended the raw tap coordinates so the line reaches the pin.

\subsection{Deployment topologies}
The current deployment is \textbf{T1}: one national graph built from the national
extract. Per-island graphs remain supported (\textbf{T2}): one process selects a
region by snapping and, for pairs split across files, stitches registered ferry
pairs (Merak--Bakauheni, Ketapang--Gilimanuk) from two legs; unlisted pairs are
rejected with \texttt{cross\_region}. A single national graph contains ferry ways
directly, so the static connector is not on its normal path.

The HTTP server exposes \texttt{GET}/\texttt{POST /route}, with
\texttt{json}, \texttt{points}, \texttt{geojson}, and \texttt{polyline} output;
\texttt{GET /healthz}, \texttt{/version}, and \texttt{/metrics}; ETag validation
before rate limiting; and structured error reasons \texttt{snap},
\texttt{too\_far}, \texttt{budget}, \texttt{cross\_region}, and
\texttt{deadline}. The deployment is one static binary (6.7\,MB) in a 36.7\,MB
container image, with the graph mounted read-only.

\section{Evaluation setup}
Numbers come from four machines and are not mixed across claims.
\begin{itemize}
  \item \emph{E1} (OSRM comparison): AMD Ryzen 5 5600G (6 cores, 12 threads),
        16\,GB RAM, Linux 7.2.15; OSRM v5.26.0 MLD in Docker; RouteID at the same
        commit. Java extract (857.6\,MB PBF), 600 origin--destination pairs, 120
        per straight-line bucket, sampled once. Requests alternate between
        engines; percentiles use successful, filter-matched queries.
  \item \emph{National build}: the same Ryzen, 12 workers, with
        \texttt{GOMEMLIMIT} set to 5\,GiB; 11\,m40\,s includes the PBF download.
  \item \emph{100\,km cap}: the same Ryzen, three Java region graphs, not the
        single national graph. This distinction matters and is repeated wherever
        the number appears.
  \item \emph{Micro-optimisations} ($\varepsilon$, ordering): an Intel Core
        i5-7360U (2 cores, 4 threads), 7\,GB RAM, on an earlier format-v3 graph.
        These are trend measurements, not final v5 latency.
  \item \emph{Throughput}: an Intel i3-6100T (2 cores, 4 threads), 7.8\,GB RAM,
        container limited to 1.5\,GiB.
\end{itemize}

The E1 filter admits a pair only when both engines snap both endpoints within
250\,m and the two snapped points lie within 50\,m of each other. Of 600 pairs,
332 enter the distance and latency statistics, 79 fail in RouteID, and 189 are
excluded by the filter. Failed queries are reported, not hidden. The filter is
necessary because a pair whose endpoints are kilometres from a road is not a
like-for-like routing comparison; it is also a limitation, discussed in
\S\ref{sec:threats}.

\section{Results}
\subsection{Artefacts and preprocessing}
\begin{table}[t]
\small
\begin{tabular}{@{}lrrrr@{}}
\toprule
graph & nodes & edges & bytes & build \\
\midrule
RouteID Java v4 (E1) & 22.71\,M & 46.68\,M & 693.1\,MB & 28\,s \\
OSRM MLD Java        & ---      & ---      & 5.08\,GB  & 250\,s \\
RouteID Java v5      & 23.46\,M & 48.32\,M & 693.3\,MB & --- \\
RouteID Sumatra v5   & 11.95\,M & 24.49\,M & 353.3\,MB & --- \\
RouteID Nusa v5      &  5.06\,M & 10.23\,M & 148.8\,MB & --- \\
RouteID Indonesia v5 & 54.89\,M & 112.34\,M & 1{,}621.2\,MB & 11\,m40\,s \\
\bottomrule
\end{tabular}
\caption{Artefacts. E1 used format v4 with the car profile; deployed v5 carries
car and motorcycle in one file, so E1 node/edge counts are not directly
comparable. The national build time includes downloading the PBF; extraction
alone is 1\,m12\,s. OSRM peak preprocessing RSS was 6.25\,GB; RouteID's national
build was capped at 5\,GiB.}
\label{tab:artefacts}
\end{table}

Table~\ref{tab:artefacts} is the first result: at the Java scale the RouteID
artefact is 7.3$\times$ smaller than OSRM's MLD data, and preprocessing is
9$\times$ faster. Part of that gap is expected---MLD buys query speed with
auxiliary data---and part is the narrower feature set: no instructions, no
alternative routes, no matrix, and no map matching.

\subsection{Latency and the OSRM comparison}
\begin{table}[t]
\footnotesize
\begin{tabular}{@{}lrrrrr@{}}
\toprule
bucket & matched & failed & excluded & RouteID p50/p95 & OSRM p50/p95 \\
       &         &        &          & (ms)            & (ms) \\
\midrule
0--5\,km   & 86 &  0 & 34 &   4.0/13.8 & 0.7/1.0 \\
5--15\,km  & 73 &  6 & 41 &  21.7/72.3 & 1.2/1.7 \\
15--30\,km & 67 & 10 & 43 &  64.8/271.1 & 1.6/2.3 \\
30--50\,km & 56 & 27 & 37 & 164.6/548.9 & 2.2/3.2 \\
50--80\,km & 50 & 36 & 34 & 525.9/1{,}675.3 & 3.2/5.1 \\
\midrule
all        & 332 & 79 & 189 & 38.7/674.0 & 1.5/3.7 \\
\bottomrule
\end{tabular}
\caption{E1 latency on the same 600 pairs, Ryzen 5 5600G, single-threaded
RouteID. Percentiles use the 332 filter-matched pairs; failures are counted per
straight-line bucket. OSRM is faster at every bucket.}
\label{tab:e1lat}
\end{table}

Table~\ref{tab:e1lat} is the central performance result. OSRM is about
$26\times$ faster at the median and $180\times$ faster at p95. RouteID is not a
low-latency replacement for MLD. Its advantage is the resource profile: 693\,MB
versus 5.08\,GB on disk, 28\,s versus 250\,s to preprocess, and 17.3\,MB versus
2.36\,GiB resident before the first query. For a workload that values a small
artefact and a restart without a multi-gigabyte warm-up, that is the relevant
trade; for a workload that values per-query latency, OSRM wins clearly.

\subsection{Distance, duration, and failures}
On the 332 matched pairs, the relative route-distance difference (RouteID minus
OSRM, divided by OSRM) has median 0.23\,\%, p95 14.49\,\%, and maximum 70.44\,\%.
Twenty-five pairs exceed 10\,\% and eight exceed 25\,\%. The pre-registered gate
was median below 0.1\,\% and p95 below 12\,\%; both conditions fail. On duration,
the signed difference has median $-1.11$\,\% and p95 $+2.39$\,\%; the p95 of the
absolute difference is 10.22\,\%. D4 therefore made absolute duration p95 the
primary metric, but this is still an engine-to-engine comparison, not accuracy.

RouteID returns 422 for 79 of the 600 pairs: 72 because an endpoint has no road
within the 3\,km snap radius, and 7 because the search budget is exhausted. There
are no \texttt{no\_route} failures in the final run. OSRM succeeds on all 79, but
its target snap is a median of 15.3\,km from the input (p90 32.7\,km, maximum
44.6\,km); those pairs are not a like-for-like comparison and are not counted as
RouteID wins. The remaining 189 pairs are excluded by the symmetric snap filter.

The residual distance spread is not a restriction defect. An audit of all 593
skipped restriction relations and the 20 largest deviations
(\texttt{paper/RESULTS.md}\,\S18) attributes it to differing cost functions and
route choice: RouteID minimises drive time and keeps \texttt{track} edges, while
OSRM minimises its \texttt{weight}. Aligning the fallback speed table with
\texttt{car.lua} reduced p95 from 21.05\,\% to 14.49\,\%, but the tail remains.

\subsection{Heuristic inflation}
\begin{table}[t]
\footnotesize
\begin{tabular}{@{}lrrrr@{}}
\toprule
$\varepsilon$ & p50 15--30\,km & p50 50--80\,km & route length & latency ratio \\
\midrule
1.0 &  87.8\,ms & 582.6\,ms & optimal & 1.00 \\
1.5 &  62.3\,ms & 387.6\,ms & same, 526/526 & 0.72 \\
2.0 &  44.7\,ms & 222.9\,ms & one pair $+11$\,\% & 0.52 \\
\bottomrule
\end{tabular}
\caption{Inflation versus optimality, paired over 526 successful queries on the
i5/format-v3 setup. Distance excess is relative to $\varepsilon=1.0$ per query.
This is a trend measurement; it was not re-run on the final v5 graph.}
\label{tab:eps}
\end{table}

Inflation is nearly free in accuracy at this radius (Table~\ref{tab:eps}). Because
the geographic bound is weak in dense networks, the optimality-preserving search
wastes most of its work; a factor of two in $\varepsilon$ buys roughly a factor of
two in median latency. We ship $\varepsilon=1.5$ because it reproduced the optimal
length on every measured pair; $\varepsilon=2.0$ is faster but introduces a
visible outlier.

\subsection{Node ordering}
Renumbering lowers p50 by 7--20\,\% and p95 by 7--16\,\% on the i5/format-v3
setup, with identical expansion counts. Its larger effect is on memory: replaying
29 metropolitan pairs raises resident set to 1{,}225\,MB with OSM-id order and
675\,MB with Hilbert order, a factor of 1.8. The gain is real but modest; early
single-corridor experiments overstated it because they touched few pages. The
result is a layout effect, not an algorithmic speed-up, and we do not extrapolate
it to server CPUs without measurement.

\subsection{Memory and workspace}
\begin{table}[t]
\footnotesize
\begin{tabular}{@{}lrrl@{}}
\toprule
state & RSS & anonymous & note \\
\midrule
RouteID E1, start        &  17.3\,MB & --- & one Java graph \\
RouteID E1, 600 pairs    & 815.3\,MB & --- & matched workload \\
OSRM MLD, before queries &  2.36\,GiB & --- & same machine \\
OSRM MLD, after 600 pairs & 2.45\,GiB & --- & same machine \\
National v5, idle        &  11.1\,MB & 4.5\,MB & header only \\
National v5, 4$\times$96\,km & 225\,MiB & 90\,MiB & concurrent \\
National v5, one 1{,}882\,km & 701\,MiB & 401\,MiB & largest tested \\
\bottomrule
\end{tabular}
\caption{Resident memory. E1 numbers are from the Ryzen comparison; national
numbers are a separate build and workload. The national process maps 1.621\,GB
but touches only the pages a query needs, which is why idle RSS is below the
artefact size.}
\label{tab:memory}
\end{table}

The workspace is a shared pool of 12-byte slots (a packed key and generation
stamp, g-score, and came-from edge). At the 100\,km cap, a 256\,MB workspace cap produced 201\,MB of slots
(16.8\,M) and the A* open set still exhausted before the workspace did; raising
the cap to 512\,MB or 2\,GB did not move the pair that failed. The pool shrinks
when idle: one run went from 10{,}485{,}760 slots to 262{,}144, with
\texttt{RssAnon} falling from 143.5\,MB to 4.5\,MB. Memory is therefore budgeted
by longest route times concurrency, and the pool is a ceiling rather than a
per-query allocation.

\subsection{Turn restrictions and snapping}
The Java build used for compliance evaluation contains 4{,}654
\texttt{type=restriction} relations; 3{,}962 remain after the profile and
reachability filter, resolving to 3{,}808 forbidden transitions and 154
permit-lists (30.1\,KB, 0.004\,\% of the 693\,MB graph; a dense CSR over edges
would need 374\,MB). We route the tail of the incoming edge to the head of the
outgoing edge for every stored restriction and check whether the returned path
still contains the forbidden transition. Without enforcement, 96.7\,\% of the
3{,}463 routable restrictions are violated, so the harness is not vacuous. With
enforcement the violation count is \textbf{0} and the fallback count is
\textbf{0}. A build-time edge-state search relaxes 116 bans whose endpoints have
no alternative path once every other ban and permit-list is enforced; the 57
searches that hit the one-million-state cap are relaxed too. The harness tests
only stored bans, so those 116 rules are outside its sample. Median latency
changes by less than a millisecond per bucket and expansion counts by 0.1\,\%.

The deployed v5 car graph stores 3{,}694 bans and 154 permit-lists; the motorcycle
profile stores 3{,}589 and 136. The reduction from 3{,}808 reflects the final
profile and reachability pass, not a relaxation of enforcement. Time-dependent
rules and via-way restrictions that require edge splitting are counted but not
applied; 593 relations are skipped for these reasons and for unresolved
references. They are a known gap, not a hidden success.

\begin{table}[t]
\small
\begin{tabular}{@{}lrrrr@{}}
\toprule
segment length & n & node p50 & node p95 & segment p50 / p95 \\
\midrule
$\le$20\,m      & 3{,}382 & 17.2\,m & 36.8\,m & 15.8 / 36.1\,m \\
20--50\,m        & 1{,}094 & 20.8\,m & 38.7\,m & 15.7 / 37.0\,m \\
50--150\,m       &   477   & 31.9\,m & 53.2\,m & 18.6 / 36.7\,m \\
$>$150\,m        &    47   & 67.3\,m & 114.0\,m & \textbf{16.1 / 34.9}\,m \\
\bottomrule
\end{tabular}
\caption{Attachment error for 5{,}000 synthetic taps pushed 0--40\,m
perpendicularly off a segment (western Sumatra, 1.02\,M segments in scope). The
injected offset is itself 0--40\,m, so the segment column measures proximity to
the road; the node column adds half a segment on top. The median segment in this
extract is only 11\,m because every OSM node is a graph node, which is why the
gain concentrates on long inter-urban segments.}
\label{tab:snap}
\end{table}

Segment snapping removes the first-order attachment error. On segments longer
than 150\,m, the median error falls from 67.3\,m to 16.1\,m
(Table~\ref{tab:snap}); on a real 22.7\,km Medan--Binjai route the tap offset drops
from 45\,m to 4\,m, both ends are trimmed, and the reported distance falls by
58\,m. The residual is the spur case: when the path leaves the anchor away from
the snapped segment, the offset becomes a straight line. That is an explicit
limitation, not an accuracy claim.

\subsection{National build and the 100\,km cap}
\begin{table}[t]
\footnotesize
\begin{tabular}{@{}lrrrr@{}}
\toprule
setup & success & p99/max & failures \\
\midrule
Ryzen, 3 regions, 100\,km & 481/600 & ---/9.5\,s & 72 snap, 39 too\_far, 8 budget \\
i3, 6\,s budget & 472/600 & 2.94/4.74\,s & 40 too\_far, 16 budget \\
i3, no budget, 30\,s & 475/600 & 4.52/9.06\,s & 45 too\_far, 8 exhausted \\
\bottomrule
\end{tabular}
\caption{100\,km-cap behaviour. The Ryzen row uses three Java region graphs,
\emph{not} the single national graph; wall-clock p50 is 39--47\,ms because short
pairs dominate, and p95 is 0.94--1.32\,s. The i3 rows use a 1.5\,GiB container.
The 6\,s budget trades three successes for a shorter tail and no exhausted
searches.}
\label{tab:cap}
\end{table}

A national graph is now a build result, not only an estimate: 54.89\,M nodes,
112.34\,M edges, 1.621\,GB, built with 12 workers under a 5\,GiB
\texttt{GOMEMLIMIT} in 11\,m40\,s including download (1\,m12\,s extraction). The
process starts in 10\,ms and idles at 11.1\,MB. What is \emph{not} measured is
query latency on that single national graph. The 100\,km-cap numbers in
Table~\ref{tab:cap} come from three Java regions on the Ryzen, and we label them
as such.

At the cap, 481 of 600 pairs succeed. Failures are 72 snap, 39 \texttt{too\_far},
and 8 budget/exhausted. The wall-clock p50 is 39--47\,ms and p95 is
0.94--1.32\,s, with a maximum of 9.5\,s. The tail is the cost of omitting a
hierarchy: a 50--80\,km pair has RouteID p95 1.68\,s in E1, and the cap run
extends that tail. On the i3, a 6\,s budget returns 472 pairs, 16 of them marked
budget; without a budget, 475 pairs succeed but 8 searches exhaust and the p99
rises from 2.94\,s to 4.52\,s. The budget is a latency control, not a free win.

\subsection{Throughput}
\begin{table}[t]
\small
\begin{tabular}{@{}lrrrr@{}}
\toprule
route class & server p50 & clients & RPS & client p95 \\
\midrule
7\,km   &  47\,ms & 4 & 75   & --- \\
15\,km  & 136\,ms & 4 & 17   & --- \\
42\,km  &  93--102\,ms & 4 & 36--40 & 105--145\,ms \\
148\,km & 4.2\,s  & 4 & 0.6  & 6.9\,s \\
mix $\le$40\,km, 298 pairs & 73.6\,ms & 4 & 23.5--24.9 & 705.4\,ms \\
\bottomrule
\end{tabular}
\caption{Throughput on the i3-6100T deployment (2 cores/4 threads, 1.5\,GiB
container, \texttt{-conc 4}), keep-alive clients. The 42\,km row is the
post-optimisation range; the mix row is a real distance distribution, not a
single class. Error rate was zero in all runs (1{,}000+ requests).}
\label{tab:throughput}
\end{table}

With four clients, a single 42\,km class sustains 36--40\,RPS at p50
93--102\,ms, while a real mix capped at 40\,km sustains 23.5--24.9\,RPS at p50
73.6\,ms and p95 705.4\,ms. The difference is queueing: the mix contains many
short routes that finish quickly but also long ones that occupy a core. On two
physical cores, concurrency above two buys no throughput. At 148\,km the server
is at 0.6\,RPS, which is another symptom of the missing hierarchy.

\section{Discussion}
\paragraph{Where the design wins.}
(i)~Footprint: 17.3\,MB resident for the E1 Java graph and 11.1\,MB idle for the
national graph, with a 10\,ms start-up. The router can share a small box with
other services and restart without a multi-gigabyte warm-up. (ii)~Operational
simplicity: one 6.7\,MB static binary, one graph file, no Lua profile at runtime,
and a monthly rebuild measured in minutes per island. (iii)~Auditability:
restrictions are compiled into the graph, reported at \texttt{/healthz}, and
checked by a harness that shows 0 violations after enforcement versus 96.7\,\%
without it. (iv)~Predictable memory: growth is bounded by longest route times
concurrency, and the workspace shrinks when idle.

\paragraph{Where it loses.}
OSRM is faster by 26$\times$ at the median and 180$\times$ at p95 on the same
data. RouteID's distance tail also fails the pre-registered gate. A product that
needs per-query latency, accurate congestion ETAs, worldwide coverage, matrix
queries, or turn-by-turn instructions should use an established engine or a
hosted API. The case for RouteID is narrow: bounded-radius car/motorcycle routes,
static speeds, a small VPS, and no per-request fee.

\paragraph{Why the tail exists.}
Without a hierarchy, the search explores a disc whose area grows with the route
radius; dense metropolitan corridors make it worse because many short edges lie
inside the disc. At 50--80\,km the E1 p95 is 1.68\,s, and at the 100\,km cap the
p95 is 0.94--1.32\,s with 8 budget failures. A contraction hierarchy or a
two-level highway/core decomposition would fix the tail but adds preprocessing
cost and auxiliary memory, which is the trade this design deliberately avoids.

\paragraph{Accuracy.}
Segment snapping was the first-order fix, not turn restrictions. The remaining
error sources are the spur case, static speeds, unapplied time-dependent and
via-way restrictions, and the choice of cost function. The distance comparison
against OSRM is useful as a regression signal but cannot establish accuracy: both
engines consume the same OSM graph, and E2 was not run. We state this in the
abstract and in \S\ref{sec:threats} rather than qualifying a claim we cannot
support.

\paragraph{Artifact and reproducibility.}
The repository contains the raw CSVs, the measurement ledger, the harness
scripts, the graph format version, PBF SHA-256 values, and the commands that
produced each table. The repository is currently private, so we do not claim a
public artefact or a reusable benchmark suite. The numbers in this draft are
reproducible from the recorded inputs on the stated machines, not from a container
image published with the paper.

\section{Threats to validity}
\label{sec:threats}
(i)~\emph{Hardware and runs.} Measurements come from a small number of commodity
machines, not server CPUs; latency is single-run, and the $\varepsilon$ and
ordering experiments use an older format-v3 graph. The E1 comparison alternates
engines but does not repeat the full run five times. Locality results in
particular should be re-measured on different cache and page sizes.

(ii)~\emph{No independent ground truth.} E2 was not run. OSRM is not a reference:
it consumes the same OSM data and has its own cost function. The E1 median
difference is 0.23\,\%, but p95 is 14.49\,\% and the maximum is 70.44\,\%, so the
pre-registered gate fails and no accuracy claim is possible. The owner chose the
limitation-only option on 2026-11-10; this paper does not contain an accuracy
result.

(iii)~\emph{Filtering and sample size.} Of 600 pairs, 332 are matched, 79 fail,
and 189 are excluded. Failed endpoints are excluded from percentiles, which
flatters latency; we report the counts and the OSRM snap distance on failures
(median 15.3\,km) so the exclusion is visible. The p99 estimates rest on 332
queries and are weak.

(iv)~\emph{Snapping study.} The attachment error uses synthetic taps with
injected offsets, not recorded user taps. The 5{,}000 samples come from western
Sumatra and may not represent dense Jakarta or rural Papua.

(v)~\emph{National scope.} The national graph is built and its memory is measured,
but its query latency is not. The 100\,km-cap numbers use three Java regions;
they are not a full-national benchmark. Ferry behaviour is static and does not
include schedules.

(vi)~\emph{Feature scope.} There is no traffic, no time-dependent restriction, no
via-way edge splitting, no alternative routes, no matrix, and no map matching.
Durations are static-speed estimates, not ETAs.

\section{Related work}
\paragraph{Shortest-path speed-up techniques.}
Dijkstra's algorithm~\cite{dijkstra1959} and A*~\cite{hart1968} are the
foundation. Goal-directed search with landmark lower bounds (ALT) was introduced
by Goldberg and Harrelson~\cite{goldberg2005}. Contraction
hierarchies~\cite{geisberger2008,geisberger2012} and customisable route
planning~\cite{delling2011,delling2017} (the MLD family used by OSRM) achieve
microsecond queries by preprocessing; their cost is preprocessing time and
auxiliary memory. Highway-dimension analyses~\cite{abraham2010,abraham2016}
explain why road networks admit such hierarchies. These techniques optimise query
time for long routes; we instead optimise resident memory for a bounded-radius
workload, and show where the simpler search stops being competitive.

\paragraph{Open-source engines.}
OSRM~\cite{osrm}, GraphHopper~\cite{graphhopper}, Valhalla~\cite{valhalla},
BRouter~\cite{brouter}, and RoutingKit~\cite{dibbelt2014} are the practical
baselines. They are general-purpose: multiple profiles, matrices, map matching,
turn-by-turn instructions. Our claim is not that we beat them on query latency;
it is that a constrained workload permits a much smaller artefact. The E1 numbers
make the trade explicit rather than hiding it behind a throughput claim.

\paragraph{Locality and out-of-core graph processing.}
Storing graphs so that traversal is page-local has a long tradition in
out-of-core graph systems (GraphChi~\cite{kyrola2012}, X-Stream~\cite{royle2013},
FlashGraph~\cite{zheng2015}) and in shared-memory frameworks such as
Ligra~\cite{shun2013}. Space-filling curves date back to
Hilbert~\cite{hilbert1891}. We apply the same idea to a routing graph: numbering
nodes along a Hilbert curve makes an A* frontier---which is spatially
compact---contiguous in the file. The measured gain is modest, and we report it
as such.

\section{Conclusion and future work}
For bounded-radius routing, the interesting resource is not query time but
resident memory, and memory-mapping plus locality-aware node numbering delivers
it with a small, auditable code base. The costs are equally clear: OSRM is much
faster, the distance tail is wider than the gate we set, and long routes need a
hierarchy we did not build. The immediate work is an independent-reference
comparison (E2), time-dependent and via-way restrictions, a full-national latency
benchmark, and a streaming builder that keeps the whole country in one file on an
8\,GB machine. A hierarchy should be added only if the product needs routes beyond
the current cap.

\bibliographystyle{ACM-Reference-Format}
\bibliography{refs}

\end{document}
